\documentclass[12pt,a4paper]{report}
\usepackage[utf8]{inputenc}
\usepackage[T1]{fontenc}
\usepackage[romanian]{babel}
\usepackage{amsmath,amsthm,amssymb,graphicx,hyperref}
\usepackage[left=1.2in,right=1in,top=1in,bottom=1in]{geometry}
\usepackage{indentfirst}
\setlength{\parindent}{1.2cm}
\setlength{\parskip}{6pt}

% ======================================
% FORMAT CAPITOLE (I. Introducere etc.)
% ======================================
\usepackage{titlesec}
\renewcommand{\thechapter}{\Roman{chapter}}
\titleformat{\chapter}
  {\normalfont\huge\bfseries}
  {\thechapter.}
  {1em}
  {}

\def\contentsname{Cuprins}
\def\bibname{Bibliografie}
\def\figurename{Figura}
\def\tablename{Tabelul}

\begin{document}

% ============================
% PAGINA DE TITLU
% ============================

\thispagestyle{empty}
\begin{center}
{\large \textbf{UNIVERSITATEA DE VEST DIN TIMIȘOARA}}\\[4pt]
{\large \textbf{FACULTATEA DE INFORMATICĂ}}\\[4pt]
{\large \textbf{PROGRAMUL DE STUDII DE LICENȚĂ: INFORMATICĂ }}\\[8em]

{\Huge \textbf{LUCRARE DE LICENȚĂ}}\\[4em]
{\LARGE \textbf{Aplicație pentru transfer de fișiere în LAN}}\\[10em]

\begin{tabular}{p{6cm}p{6cm}}
\textbf{COORDONATOR:} & \textbf{ABSOLVENT:}\\
Asist. Dr. Laurențiu Coroban & Capac Ștefan--Anelin
\end{tabular}

\vfill
{\large TIMIȘOARA, 2026}
\end{center}

\newpage

% ============================
% ABSTRACT
% ============================

\section*{Abstract}
\addcontentsline{toc}{section}{Abstract}

Prezenta lucrare de licență propune \textit{Fishare}, o aplicație desktop orientată spre transferul rapid și sigur de fișiere între dispozitive aflate în aceeași rețea locală (LAN). Soluția urmărește reducerea dependenței de servicii cloud și oferirea unui mecanism direct de comunicare între utilizatori, cu accent pe confidențialitate, control local și experiență simplă de utilizare.

Aplicația este construită pe o arhitectură hibridă: interfața grafică și logica de control sunt realizate în Python (PyQt6), iar transferul efectiv al datelor este susținut de module C++, pentru o experiență fluentă. În plus, sistemul include descoperire automată a dispozitivelor din rețea, alegerea inteligentă a protocolului de transfer și evidență clară a operațiilor efectuate.

Rezultatul este o platformă practică pentru scenarii locale (laboratoare, birouri, rețele casnice), care combină performanța cu măsuri de securitate moderne și cu un flux de utilizare intuitiv.

\newpage

% ============================
\tableofcontents
\newpage

% ============================
\chapter{Introducere}
% ============================

Într-un context în care transferul de date rapide și sigure reprezintă o componentă esențială a colaborării digitale, soluțiile tradiționale de partajare a fișierelor -- bazate pe servicii online sau pe aplicații de mesagerie -- devin tot mai limitate. Acestea depind de conexiunea la internet, pot compromite confidențialitatea datelor și implică pași tehnici suplimentari, care pot fi dificili pentru utilizatorii non-tehnici.

Lucrarea de față își propune să dezvolte o aplicație software modernă, denumită \textbf{Fishare}, care să faciliteze transferul de fișiere între dispozitive conectate la aceeași rețea locală (LAN), fără a necesita servere intermediare sau configurări complexe.

\section*{Motivația lucrării}

Nevoia unei aplicații dedicate pentru transferul local de fișiere provine din lipsa unei soluții rapide și prietenoase care să funcționeze independent de internet. În multe medii -- laboratoare universitare, birouri, rețele de companie -- utilizatorii colaborează în aceeași rețea locală și au nevoie de o metodă de schimb de fișiere intuitivă, sigură și instantanee.

Aplicațiile existente, precum serviciile cloud, implică stocarea fișierelor pe servere externe, ridicând probleme legate de securitate și confidențialitate. Alte metode, precum partajarea directoarelor de rețea, sunt lente, greu de configurat și vulnerabile. Din acest motiv, Fishare a fost conceput pentru a elimina aceste neajunsuri și pentru a permite o experiență directă, locală și controlată de utilizator.

\section*{Definirea problemei}

Problema identificată este lipsa unei aplicații simple, performante și autonome care să permită trimiterea fișierelor între calculatoare dintr-o rețea locală, cu detecție automată a dispozitivelor și mecanisme de control al vizibilității.

O soluție optimă trebuie să asigure:
\begin{itemize}
    \item transferuri rapide și fiabile;
    \item securitatea datelor trimise;
    \item confidențialitate totală, fără servere intermediare în cloud;
    \item handshake între peers pentru permiterea transferului;
    \item o interfață grafică intuitivă;
    \item gestionarea stării (Disponibil / Ocupat);
    \item opțiuni de control pentru notificări și solicitări de transfer.
\end{itemize}

\section*{Soluția propusă}

Aplicația oferă o soluție completă, bazată pe descoperirea automată a dispozitivelor din rețea și pe schimbul direct de fișiere. Interfața grafică este realizată în PyQt, iar transferul este gestionat prin QUIC, cu fallback la TCP.

Aplicația permite trimiterea simultană a mai multor fișiere, configurarea directorului de descărcare și schimbarea numelui vizibil. Fluxul este gândit pentru utilizare simplă: selectare fișiere, selectare dispozitiv(e), confirmare și monitorizare progres.

\begin{figure}[!ht]
    \centering
    \includegraphics[width=0.9\linewidth]{fishrep.png}
    \caption{Reprezentare grafică a funcționalității aplicației \textit{Fishare}}
\end{figure}

\newpage

% ============================
\chapter{State of the Art}
% ============================

Transferul de fișiere între dispozitive poate fi realizat prin mai multe categorii de soluții. Prima categorie include platformele cloud, care oferă acces facil și sincronizare, dar depind de internet și presupun stocare externă. A doua categorie include partajarea clasică în rețea locală (de tip foldere partajate), utilă în medii controlate, dar adesea dificil de configurat pentru utilizatorii obișnuiți. A treia categorie include aplicațiile peer-to-peer locale, care oferă transfer direct și latență redusă.

Tendința actuală în aplicațiile de tip LAN sharing este orientarea către trei direcții: simplitatea fluxului de utilizare, securizarea end-to-end și folosirea de mecanisme concurente pentru performanță mai bună. În acest context, Fishare adoptă o abordare local-first: descoperire automată în LAN, conexiune directă între dispozitive și procesare eficientă a datelor transferate.

Comparativ cu abordările tradiționale, soluția propusă în această lucrare urmărește un echilibru între ușurința de utilizare și controlul tehnic asupra confidențialității și performanței.

\newpage
\chapter{Modelare și Design (UML)}

Acest capitol prezintă modelele UML utilizate pentru a descrie structura și comportamentul aplicației. Prin aceste diagrame sunt evidențiate funcționalitățile principale, interacțiunile dintre utilizator și sistem, precum și fluxurile interne de procesare, oferind o imagine clară asupra modului în care aplicația este organizată și funcționează.

\section*{Diagrama de tip Use Case}

\begin{center}
    \includegraphics[width=0.85\linewidth]{usecase.png}
    \\[0.3em]
    \small \textit{Diagramă de tip Use Case pentru aplicația Fishare}
\end{center}

\section*{Descrierea principalelor Use Case-uri}

\subsection*{Selectarea fișierelor}
Utilizatorul poate adăuga unul sau mai multe fișiere într-o listă de trimitere. Lista poate fi modificată înainte de inițierea transferului, iar fișierele pot fi eliminate individual sau în bloc.

\subsection*{Selectarea dispozitivelor (peers)}
Aplicația afișează automat dispozitivele detectate în LAN. Utilizatorul alege unul sau mai multe dispozitive țintă, în funcție de necesitățile curente.

\subsection*{Inițierea transferului}
După selecție, utilizatorul inițiază trimiterea. Sistemul creează cererea de transfer, negociază protocolul disponibil și pornește transferul în fundal.

\subsection*{Descoperirea automată a dispozitivelor}
Dispozitivele active sunt identificate prin anunțuri periodice în rețea locală. Lista este actualizată continuu, iar dispozitivele inactive sunt eliminate automat după expirarea unui interval de timp.

\subsection*{Gestionarea transferului de fișiere}
Aplicația gestionează progresul transferului pentru fiecare dispozitiv și actualizează indicatorii din interfață. În cazul erorilor, sunt salvate informații în istoric și se aplică politici de curățare pentru fișierele incomplete.

\subsection*{Acceptarea sau respingerea unui transfer primit}
La primirea unei solicitări, utilizatorul primește un dialog de confirmare. Transferul este acceptat sau respins explicit, iar la lipsa răspunsului într-un interval prestabilit solicitarea este tratată ca refuz.

\subsection*{Configurarea aplicației}
Utilizatorul poate configura numele dispozitivului, folderul de descărcare și starea de disponibilitate. Setările sunt persistate local pentru rulările viitoare.

\newpage
\chapter{Arhitectura și tehnologii utilizate}

\section*{Arhitectura aplicației}

Arhitectura Fishare este modulară și separă clar zonele de responsabilitate:
\begin{itemize}
    \item \textbf{Interfață grafică (PyQt6):} afișează dispozitive, fișiere selectate, progres și ferestre de confirmare;
    \item \textbf{Stare aplicație:} menține thread-safe dispozitivele active, statusul local și progresul transferurilor;
    \item \textbf{Discovery în LAN:} advertiser și scanner pentru detectarea automată a peer-ilor;
    \item \textbf{Serviciu de transfer:} gestionează cozi per dispozitiv, paralelismul și logica de retry;
    \item \textbf{Protocoale de comunicație:} QUIC (preferat) și TCP (fallback);
    \item \textbf{Securitate:} schimb de chei, autentificare peer și criptare AEAD;
    \item \textbf{Persistență locală:} configurație, dispozitive cunoscute și istoric transferuri.
\end{itemize}

\section*{Tehnologii utilizate}

\begin{itemize}
    \item \textbf{Python 3} -- orchestrare, UI și logică aplicație;
    \item \textbf{PyQt6} -- dezvoltarea interfeței grafice desktop;
    \item \textbf{C++17 + pybind11} -- accelerarea operațiilor de transfer;
    \item \textbf{OpenSSL / cryptography} -- primitive criptografice și AEAD;
    \item \textbf{aioquic} -- suport pentru protocolul QUIC (opțional);
    \item \textbf{JSON} -- format de stocare pentru date persistente.
\end{itemize}

\newpage
\chapter{Implementare și funcționalități}

\section*{Facilități ale aplicației}

În forma actuală, aplicația oferă următoarele facilități principale:
\begin{itemize}
    \item \textbf{Conectare imediată cu dispozitivele din jur:} aplicația detectează automat dispozitivele active din aceeași rețea locală;
    \item \textbf{Trimitere simplă către unul sau mai mulți destinatari:} utilizatorul selectează fișierele și alege rapid cui dorește să le trimită;
    \item \textbf{Feedback vizual în timp real:} progresul și viteza transferului sunt afișate clar, astfel încât utilizatorul știe permanent ce se întâmplă;
    \item \textbf{Control asupra transferurilor primite:} fiecare solicitare poate fi acceptată sau respinsă, în funcție de context;
    \item \textbf{Personalizare de bază:} numele dispozitivului și folderul de descărcare pot fi configurate ușor;
    \item \textbf{Control al disponibilității:} utilizatorul poate trece rapid între stările \textit{Available} și \textit{Busy};
    \item \textbf{Memorie a activității:} istoricul păstrează evidența transferurilor trimise și primite;
    \item \textbf{Continuitate în funcționare:} dacă un protocol nu este disponibil, aplicația poate continua transferul cu alternativa potrivită.
\end{itemize}

\section*{Cum funcționează aplicația în practică}

Aplicația este organizată astfel încât utilizatorul să poată continua interacțiunea cu interfața chiar și în timpul transferurilor. În fundal, operațiile sunt coordonate separat, astfel încât trimiterea către dispozitive diferite să se poată desfășura în paralel, fără a bloca fluxul principal.

Înainte de începerea unui transfer, aplicația verifică legătura dintre dispozitive și pregătește comunicarea într-un mod sigur. Datele sunt protejate pe durata transmiterii, iar fișierele primite sunt salvate controlat, în locația setată de utilizator.

Pentru a menține o experiență rapidă și stabilă, componentele cele mai solicitante sunt gestionate de un motor C++ integrat în aplicația Python. Această combinație oferă un echilibru bun între flexibilitate la dezvoltare și performanță în utilizarea de zi cu zi.

\newpage
\chapter{Testare și evaluarea performanței}

\newpage
\chapter{Concluzii și direcții de dezvoltare}

\newpage
\chapter{Plan de activitate}

\section*{Etape propuse}
\begin{itemize}
    \item \textbf{Octombrie 2025:} Analiză și documentare a protocoalelor de comunicație (QUIC, TCP) și investigarea posibilităților de implementare în tehnologiile utilizate în proiect.

    \item \textbf{Noiembrie 2025:} Studierea soluțiilor similare existente și definirea obiectivelor de dezvoltare.

    \item \textbf{Decembrie 2025:} Proiectarea arhitecturii generale a aplicației și definirea principalelor cazuri de utilizare.

    \item \textbf{Ianuarie 2026:} Dezvoltarea interfeței grafice în PyQt și implementarea componentelor de configurare a aplicației.

    \item \textbf{Februarie -- Martie 2026:} Implementarea logicii de transfer și a protocoalelor de comunicație în C++ și integrarea modulelor cu partea de aplicație dezvoltată în Python.

    \item \textbf{Aprilie 2026:} Testarea, validarea și optimizarea performanței aplicației.

    \item \textbf{Mai 2026:} Finalizarea documentației tehnice și pregătirea susținerii lucrării.
\end{itemize}

\begin{figure}[!ht]
    \centering
    \includegraphics[width=0.9\linewidth]{planificare.png}
    \caption{Diagrama Gantt -- Etapele de dezvoltare ale lucrării \textit{Fishare}}
\end{figure}

\newpage
\begin{thebibliography}{9}

\end{thebibliography}

\end{document}
